% GENERATED FILE. Edit canonical lessons, then run npm run generate:books.
% canonical-source-hash: fnv1a64:e8f1f1a21cc9e883
% canonical-lessons: FR-C469-humeur, FR-C469-serein, FR-C469-lunatique, FR-C469-enerver, FR-C469-agacer

\chapter{De quelle humeur ?}
\label{ch:fr-de-quelle-humeur}
\hlchaptermodality{469}

\begin{chapteropening}
\textbf{By the end of this chapter:} \emph{I can describe moods: calm, moody or irritated.}

Complete the last lesson of chapter 469: I can describe moods: calm, moody or irritated.
\end{chapteropening}

\section[l'humeur]{l'humeur — mood}

\label{lesson:FR-C469-humeur}



\begin{quote}
Before the new one: say the French for cunning, then the French for to lie.
\end{quote}

\subsection*{You'll want to know: l'humeur}
\textbf{l'humeur} — \textquotedblleft{}mood\textquotedblright{}.

\begin{quote}
Ce matin, je suis de bonne humeur.
\end{quote}

\begin{quote}
\emph{This morning I am in a good mood.}
\end{quote}

\emph{De bonne humeur} — in a good mood; \emph{de mauvaise humeur} — in a bad mood. It is feminine: \emph{une humeur}.

\begin{cousinweb}
\emph{l'humeur} is inherited from Latin \emph{ūmor} \textquotedblleft{}moisture, fluid\textquotedblright{} (EtymDB; Lewis \& Short): old medicine linked a person's mood to the fluids of the body.
\end{cousinweb}

\begin{grammarlens}[title={What you've built}]
One more word for telling, explaining and describing at B1.
\end{grammarlens}

\subsection*{Your turn}
\begin{itemize}
\raggedright
  \item \emph{Say it:} \emph{l'humeur}
  \item \emph{Say it:} \emph{l'humeur}, once more
  \item \emph{Say it:} a sentence of your own with \emph{l'humeur}
  \item \emph{Recall:} say the French for cunning, then the French for to lie, then say \emph{l'humeur} again
\end{itemize}

\begin{tcolorbox}[breakable,colback=teal!4,colframe=teal!35!black,title={Before you move on}]
What does l'humeur mean? (\textquotedblleft{}Mood\textquotedblright{}.) Say it once more.
\end{tcolorbox}
\section[serein]{serein — calm, serene}

\label{lesson:FR-C469-serein}



\begin{quote}
Before the new one: say the French for to lie, then the French for mood.
\end{quote}

\subsection*{You'll want to know: serein}
\textbf{serein} — \textquotedblleft{}calm, serene\textquotedblright{}.

\begin{quote}
Après les vacances, je me sens plus serein.
\end{quote}

\begin{quote}
\emph{After the holidays I feel calmer.}
\end{quote}

Feminine \emph{sereine}. It describes a calm mind, not loud or worried.

\begin{cousinweb}
\emph{serein} is inherited from Latin \emph{serēnus} \textquotedblleft{}clear, bright, calm\textquotedblright{} (Lewis \& Short), said first of the sky. Old French also had a noun \emph{serein} \textquotedblleft{}evening\textquotedblright{}, from Latin \emph{sērum} (EtymDB).
\end{cousinweb}

\begin{grammarlens}[title={What you've built}]
One more word for telling, explaining and describing at B1.
\end{grammarlens}

\subsection*{Your turn}
\begin{itemize}
\raggedright
  \item \emph{Say it:} \emph{serein}
  \item \emph{Say it:} \emph{serein}, once more
  \item \emph{Say it:} a sentence of your own with \emph{serein}
  \item \emph{Recall:} say the French for to lie, then the French for mood, then say \emph{serein} again
\end{itemize}

\begin{tcolorbox}[breakable,colback=teal!4,colframe=teal!35!black,title={Before you move on}]
What does serein mean? (\textquotedblleft{}Calm, serene\textquotedblright{}.) Say it once more.
\end{tcolorbox}
\section[lunatique]{lunatique — moody}

\label{lesson:FR-C469-lunatique}



\begin{quote}
Before the new one: say the French for mood, then the French for calm.
\end{quote}

\subsection*{You'll want to know: lunatique}
\textbf{lunatique} — \textquotedblleft{}moody\textquotedblright{}.

\begin{quote}
Mon frère est lunatique : un jour il rit, un jour il boude.
\end{quote}

\begin{quote}
\emph{My brother is moody: one day he laughs, the next he sulks.}
\end{quote}

Not like English \emph{lunatic}: it means moody, full of ups and downs.

\begin{cousinweb}
\emph{lunatique} is a learned borrowing from Late Latin \emph{lūnāticus} \textquotedblleft{}moonstruck\textquotedblright{} (Lewis \& Short), built on \emph{lūna} \textquotedblleft{}the moon\textquotedblright{}: moods were thought to follow its phases.
\end{cousinweb}

\begin{grammarlens}[title={What you've built}]
One more word for telling, explaining and describing at B1.
\end{grammarlens}

\subsection*{Your turn}
\begin{itemize}
\raggedright
  \item \emph{Say it:} \emph{lunatique}
  \item \emph{Say it:} \emph{lunatique}, once more
  \item \emph{Say it:} a sentence of your own with \emph{lunatique}
  \item \emph{Recall:} say the French for mood, then the French for calm, then say \emph{lunatique} again
\end{itemize}

\begin{tcolorbox}[breakable,colback=teal!4,colframe=teal!35!black,title={Before you move on}]
What does lunatique mean? (\textquotedblleft{}Moody\textquotedblright{}.) Say it once more.
\end{tcolorbox}
\section[énerver]{énerver — to annoy, to get on someone's nerves; (s'énerver) to get worked up}

\label{lesson:FR-C469-enerver}



\begin{quote}
Before the new one: say the French for calm, then the French for moody.
\end{quote}

\subsection*{You'll want to know: énerver}
\textbf{énerver} — \textquotedblleft{}to annoy, to get on someone's nerves; (s'énerver) to get worked up\textquotedblright{}.

\begin{quote}
Je m'énerve quand le train est en retard.
\end{quote}

\begin{quote}
\emph{I get worked up when the train is late.}
\end{quote}

Often reflexive: \emph{s'énerver}, past \emph{je me suis énervé}. \emph{Ça m'énerve} — that gets on my nerves.

\begin{cousinweb}
\emph{énerver} is a learned borrowing from Latin \emph{ēnervāre} \textquotedblleft{}to cut the sinews, to weaken\textquotedblright{} (EtymDB), built on \emph{nervus} \textquotedblleft{}a sinew\textquotedblright{} (Lewis \& Short). French later gave it the sense \textquotedblleft{}to set the nerves on edge\textquotedblright{}.
\end{cousinweb}

\begin{grammarlens}[title={What you've built}]
One more word for telling, explaining and describing at B1.
\end{grammarlens}

\subsection*{Your turn}
\begin{itemize}
\raggedright
  \item \emph{Say it:} \emph{énerver}
  \item \emph{Say it:} \emph{énerver}, once more
  \item \emph{Say it:} a sentence of your own with \emph{énerver}
  \item \emph{Recall:} say the French for calm, then the French for moody, then say \emph{énerver} again
\end{itemize}

\begin{tcolorbox}[breakable,colback=teal!4,colframe=teal!35!black,title={Before you move on}]
What does énerver mean? (\textquotedblleft{}To annoy, to get on someone's nerves; (s'énerver) to get worked up\textquotedblright{}.) Say it once more.
\end{tcolorbox}
\section[agacer]{agacer — to irritate, to get on someone's nerves}

\label{lesson:FR-C469-agacer}



\begin{quote}
Before the new one: say the French for moody, then the French for to annoy.
\end{quote}

\subsection*{You'll want to know: agacer}
\textbf{agacer} — \textquotedblleft{}to irritate, to get on someone's nerves\textquotedblright{}.

\begin{quote}
Le bruit de la rue agace mon voisin.
\end{quote}

\begin{quote}
\emph{The noise from the street irritates my neighbour.}
\end{quote}

Spelling: \emph{nous agaçons}, with \emph{ç} before \emph{o}.

\begin{cousinweb}
\emph{agacer} goes back to Old French \emph{agacier}. The origin is not settled: a link with \emph{agace}, an old word for the magpie, is often suggested.
\end{cousinweb}

\begin{grammarlens}[title={What you've built}]
One more word for telling, explaining and describing at B1.
\end{grammarlens}

\subsection*{Your turn}
\begin{itemize}
\raggedright
  \item \emph{Say it:} \emph{agacer}
  \item \emph{Say it:} \emph{agacer}, once more
  \item \emph{Say it:} a sentence of your own with \emph{agacer}
  \item \emph{Recall:} say the French for moody, then the French for to annoy, then say \emph{agacer} again
\end{itemize}

\begin{tcolorbox}[breakable,colback=teal!4,colframe=teal!35!black,title={Before you move on}]
What does agacer mean? (\textquotedblleft{}To irritate, to get on someone's nerves\textquotedblright{}.) Say it once more.
\end{tcolorbox}
